\chapter{参考文献格式说明}
\label{cha:chap3}


\section{参考文献格式说明}

\textbf{命令：}\textbackslash cite\{xxx\}，如深度学习\cite{lecunDeepLearning2015}，中文文献\cite{zhou2021}。


\textbf{参考文献检查项}：论文完成后应仔细检查参考文献格式，下面总结了一些在论文审阅过程中常被指出的问题：

\begin{itemize}
	\item 参考文献格式参照节\ref{subsec:ref}
	\item 除Arxiv文章，所有文献必须有页码、年份，期刊论文还应有卷号、期号
	\item 检查同一会议/期刊的名称是否统一
	\item 检查文献标题单词首字母大写是否统一
	\item 检查是否有重复文献
\end{itemize}








\subsection{参考文献格式}\label{subsec:ref}

参考文献是文中引用的有具体文字来源的文献集合。按照GB 7714《文后参考文献著录规则》的规定执行。参考文献的格式为:

著作：[序号]作者.译者.书名[M].版本(第一版不著录).出版地:出版社,出版时间:引用部分起止页.

期刊：[序号]作者.译者.文章题目[J].期刊名,年份,卷号(期数):引用部分起止页.

会议论文集：[序号]作者.译者.文章名[C]. //编者.论文集名,会议地址，会议时间.出版地:出版者，出版年.引用部分起止页.

学位论文：[序号]作者.题名[D].保存地点:保存单位,年份.引用部分起止页.

专利：[序号]专利申请者.专利文献题名[P].国别,专利文献种类,专利号.发布日期:引用部分起止页.

技术标准：[序号]起草责任者.标准代号.标准顺序号——发布年.标准名称.出版地.出版者.出版年份:引用部分起止页.

报纸: [序号]作者.题名[N].报纸名，出版日期(版次)


